% LaTeX companion of 04-linear-transformations.typ. Both formats are editable.
\chapter{Linear transformations}

\section{WS 4, p. 1: coordinate encoding and standard matrices}

现在你的位置为 \(5^{\circ}\text{E},42^{\circ}\text{N}\)。用一个 vector \(\begin{pmatrix}
5 \\
42
\end{pmatrix} \in \mathbb{R}^{2}\) 表示方位。现在你使用一个 encode 来加密你的方位：

\[\left\{ \begin{matrix}
{y_{1} = x_{1} + 3x_{2}} \\
{y_{2} = 2x_{1} + 5x_{2}}
\end{matrix} \right.,\quad\begin{pmatrix}
y_{1} \\
y_{2}
\end{pmatrix} = \begin{pmatrix}
1 & 3 \\
2 & 5
\end{pmatrix}\begin{pmatrix}
x_{1} \\
x_{2}
\end{pmatrix}.\]

这个加密方位的码为 \(\begin{pmatrix}
131 \\
220
\end{pmatrix}\). 我们做了一个 transformation，使得 the same \(\left( {y_{1},y_{2}} \right)\) but 坐标不变。例如 \(\left( {0,1} \right)\) 从 \(\left( {x_{1},x_{2}} \right)\) 坐标 map 到 \(\left( {y_{1},y_{2}} \right)\) 坐标后仍是 \(\left( {0,1} \right)\)；但 \(\left( {y_{1},y_{2}} \right)\) 下的 \(\left( {3,5} \right)\) 同样在 \(\left( {x_{1},x_{2}} \right)\) 下为 \(\left( {3,5} \right)\)，而 \(\left( {y_{1},y_{2}} \right)\) 下的 \(\left( {5,42} \right)\) 在 \(\left( {x_{1},x_{2}} \right)\) 下为 \(\left( {131,220} \right)\)。The source sketch marks these corresponding points and axes.

\begin{definition}{\kn{Linear transformation}}

A function \(T:\mathbb{R}^{m}\rightarrow\mathbb{R}^{n}\) is called a linear transformation if for every \(x \in \mathbb{R}^{m}\) there is an \(n \times m\) matrix \(A\) such that \(T(x) = Ax\).

Example: \(y = x_{1}^{2} + x_{2}^{2} + x_{3}^{2}\) has input \(\begin{pmatrix}
x_{1} \\
x_{2} \\
x_{3}
\end{pmatrix}\) and output \(\lbrack y\rbrack\); \(A = \begin{pmatrix}
x_{1} & x_{2} & x_{3}
\end{pmatrix}\). 可以发现 \(y = x \cdot x\) is not a linear transformation of \(x\).

\end{definition}

\begin{definition}{Identity matrix}

Identity matrix, denoted by \(I_{n}\): \(I_{2} = \begin{pmatrix}
1 & 0 \\
0 & 1
\end{pmatrix}\) and \(I_{3} = \begin{pmatrix}
1 & 0 & 0 \\
0 & 1 & 0 \\
0 & 0 & 1
\end{pmatrix}\).

For \(T(x) = \begin{pmatrix}
0 & {- 1} \\
1 & 0
\end{pmatrix}x\), \(T\left( \begin{pmatrix}
1 \\
1
\end{pmatrix} \right) = \begin{pmatrix}
{- 1} \\
1
\end{pmatrix}\) and \(T\left( \begin{pmatrix}
0 \\
{- 2}
\end{pmatrix} \right) = \begin{pmatrix}
2 \\
0
\end{pmatrix}\); it is a counterclockwise rotation by \(90\) degrees. Also \(x\) and \(T(x)\) have the same length: \(\sqrt{x_{1}^{2} + x_{2}^{2}} = \sqrt{\left( {- x_{2}} \right)^{2} + x_{1}^{2}}\). The source diagram labels the two arrows \(x\) and \(T(x)\).

\end{definition}

For \(T(x) = Ax\) with \(A = \begin{pmatrix}
1 & 2 & 3 \\
4 & 5 & 6 \\
7 & 8 & 9
\end{pmatrix}\), \(T\left( \begin{pmatrix}
1 \\
0 \\
0
\end{pmatrix} \right) = \begin{pmatrix}
1 \\
4 \\
7
\end{pmatrix}\) and \(T\left( \begin{pmatrix}
0 \\
1 \\
0
\end{pmatrix} \right) = \begin{pmatrix}
2 \\
5 \\
8
\end{pmatrix}\).

\section{WS 4, p. 2: linearity, bases, and transition matrices}

\begin{theorem}{Standard matrix}

Consider a linear transformation \(T:\mathbb{R}^{m}\rightarrow\mathbb{R}^{n}\). Let \(e_{i}\) be the vector whose \(i\)th component is \(1\) and all other components are \(0\). Then \(A = \begin{pmatrix}
{T\left( e_{1} \right)} & {T\left( e_{2} \right)} & \ldots & {T\left( e_{m} \right)}
\end{pmatrix}\). Equivalently, \(e_{1},e_{2},\ldots,e_{m}\) is the standard vector of \(\mathbb{R}^{n}\) (而 \(\mathbb{R}^{3}\) 中的 \(e_{1},e_{2},e_{3}\) 都用 \(i,j,k\) 来 denote).

\end{theorem}

\begin{theorem}{Linearity test}

A transformation \(T:\mathbb{R}^{m}\rightarrow\mathbb{R}^{n}\) is linear iff (a) \(\forall v,w \in \mathbb{R}^{m}\), \(T\left( {v + w} \right) = T(v) + T(w)\); and (b) \(\forall v \in \mathbb{R}^{m}\) and scalar \(k\), \(T\left( {kv} \right) = kT(v)\).

\end{theorem}

\begin{definition}{Distribution vectors and transition matrices}

A vector \(x \in \mathbb{R}^{n}\) is said to be a distribution vector if its components (1) sum to \(1\) and (2) are all \(\geq 0\). A square matrix \(A\) is a transition matrix if its col vectors are distribution vectors.

\end{definition}
